\section{Introduction}

\subsection{Motivation}
\label{sec:motivation}

The Nepal Stock Exchange (NEPSE) is a small, structurally thin emerging
market with 372 actively traded equities, thirteen official sector
taxonomies, and a combined capitalisation that is dominated by commercial
banks and hydro-power issuers.  Despite three decades of trading
history---our dataset begins 6 June 2003 and continues to
24 April 2026---research-grade quantitative analysis of
NEPSE remains sparse relative to mature Asian peers (BSE, NSE, SET, IDX).
Most published work on the exchange relies on sectoral decomposition as
the primary organising principle for diversification and risk analysis.
Yet, as the broader financial-network literature has long argued,
official sector labels are an administrative artefact that captures only
a fraction of the true covariance structure of
returns~\cite{bonanno2001topology, mantegna1999hierarchical, onnela2003dynamics,
podosu2013network}.

A long line of research---originating with Mantegna's seminal work on
hierarchical structure in financial
markets~\cite{mantegna1999hierarchical} and extended through Mantegna and
Stanley's \emph{An Introduction to
Econophysics}~\cite{mantegna1999introduction}---shows that the
\emph{correlation matrix of returns} encodes a richly organised structure
that is not visible in sector labels.  Pearson correlations can be lifted
to a metric space via $d_{ij}=\sqrt{2(1-\rho_{ij})}$, the
\emph{Mantegna distance}~\cite{mantegna1999hierarchical,onnela2003dynamics},
which is provably Euclidean and admits a Minimum Spanning Tree (MST)
sparsification~\cite{kruskal1956shortest} that retains the market's
skeletal topology while discarding the noise of weakly correlated pairs.
The MST is a particularly attractive object because, with $N-1$ edges
over $N$ nodes, it is the unique connected, acyclic, and dense-enough
summary of pairwise similarities that is suitable for community
detection on a small graph.

\subsection{Contributions}
\label{sec:contributions}

This paper accompanies and documents \textbf{NepGraph}, a fully
reproducible analytical framework that ports the correlation-network
methodology of Mantegna, Onnela and Tumminello to the NEPSE universe,
validates the sparsification parameter empirically, and exposes the
resulting structure through an interactive dashboard.

Our contributions are four-fold:
\begin{enumerate}
  \item \textbf{Pipeline.} We describe a complete, open-source pipeline
  from raw daily NEPSE close prices to an interactive minimum-spanning-tree
  visualisation, Louvain community partition, anomaly report, portfolio
  diversification score and equal-weight backtest.

  \item \textbf{Empirical parameter selection.} We perform a
  pre-registered $k$-sweep (Section~\ref{sec:ksweep}) over
  \num{440} parameter combinations to justify the choice of $k=5$
  neighbours per node in the top-$k$ graph.  Although $k=3$ slightly
  outperforms $k=5$ on the headline modularity metric, it fragments the
  network into disconnected islands, which means global topology
  metrics like betweenness centrality break down.  We argue that a
  single, connected backbone is a prerequisite for global analysis, and
  we therefore adopt the minimum $k$ that delivers it: $k=5$.

  \item \textbf{Sector-mismatch anomaly score.} We define an anomaly as
  a stock whose top-$k$ MST neighbours are dominated by a sector
  different from its official NEPSE classification, and demonstrate that
  this yields interpretable cross-sector signals (e.g.\ hydro-power
  equities co-moving with insurance issuers) that are invisible to the
  official taxonomy.

  \item \textbf{Open release.} The full code, raw $k$-sweep metrics,
  pre-rendered figures, NEPSE sector map, and analysis scripts are
  released under MIT at \url{https://github.com/sthaarwin/Nepgraph}.
\end{enumerate}

\subsection{Outline}
The remainder of the paper is organised as follows.
Section~\ref{sec:related} situates the work in the financial-network
literature.  Section~\ref{sec:method} formalises the methodology.
Section~\ref{sec:ksweep} presents the empirical $k$-sweep that calibrates
the sparsification parameter.  Section~\ref{sec:results} reports headline
results on the full NEPSE universe.
Section~\ref{sec:discussion} discusses limitations, applications and
threats to validity.  Section~\ref{sec:conclusion} concludes.

\newpage
